
\section{The Field Corpus}
\label{sec:appendix:corpus}

\subsection{Corpus construction}
\label{sec:appendix:construction}

We re-derived every version of every file rather than reuse the frame's text. A snapshot supports
claims about files; only a history supports claims about instructions. The frame is the repository
list of \citet{chatlatanagulchai2025agentreadmes}, the study that established file-level growth:
public GitHub repositories carrying at least one \texttt{CLAUDE.md}, \texttt{AGENTS.md}, or
\texttt{copilot-instructions.md} at their default branch. We cloned each bloblessly and walked the
full commit history of every tracked context file. \FieldRepos{} repositories yielded
\FieldMultiVersionFiles{} files with at least two versions, \FieldMatches{} tracked
version-to-version transitions, and \FieldSpells{} spells.

Walking forward also keeps the measurement symmetric in time. If we had anchored each file's
instruction set on its final version and searched backwards, we would have conditioned on survival
to that version. Survival is the quantity under study.

\subsection{Instruction segmentation}
\label{sec:appendix:matching}

\begin{table*}[t]
  \centering\small
  \begin{tabularx}{\textwidth}{@{}c>{\raggedright\arraybackslash}Xrrc@{}}
    \toprule
    Gate & Criterion & Threshold & Observed & Pass \\
    \midrule
    0 & Median per-repo net $\Delta D > 0$ & $> 0$ & $+7$ & $\checkmark$ \\
    0 & Grow:shrink ratio & $> 1.5$ & 2.42 & $\checkmark$ \\
    A & Deletions excluding rewrites & $\ge 300$ & 28,426 & $\checkmark$ \\
    A & Tail share at/above median age & $\ge 20\%$ & 57.0\% & $\checkmark$ \\
    B & Matcher precision / recall & $\ge 0.85$ & 1.000 / 0.933 & $\checkmark$ \\
    -- & Median rel.\ growth, all multi-version files: count vs.\ length & count $>$ length & $1.193$ vs.\ $1.028$ & $\checkmark$ \\
    \bottomrule
  \end{tabularx}
  \caption{
    Every pre-registered criterion passes on the canonical corpus. Rows: the criteria fixed before
    any pipeline code (\S\ref{sec:appendix:matching}--\ref{sec:appendix:censoring}); columns: the
    threshold each was written against and the value this run produced. Gate~B is measured on a
    50-transition hand-annotation pool rather than on this run, the only inherited number in the
    corpus half. \emph{Grow:shrink} is the share of multi-version repositories gaining instructions
    over the share losing them; \emph{tail share} is the fraction of deletions at or above the
    median observed instruction age, which guards identification past that age.
  }
  \label{tab:gates}
\end{table*}


We call one clause-level imperative an \emph{instruction} and everything else in the file
\emph{payload}. We segment in two steps: we split each file version on markdown structure, taking
list items and block boundaries, then split any remaining prose at sentence boundaries. We send
headings, fenced code, tables, whole-line bold labels, and any fragment shorter than
\SegMinWords{} words to payload, because at that length a line is usually a label and not a rule.
Payload never enters the risk set. We still count it, though, and report it beside instruction
count in \autoref{sec:growth}, so a reader can check whether the growth we report comes from
routing more text into the instruction class as files age.

We fixed the grammar per corpus, before reading the gates, and it remains our largest untested
degree of freedom on the corpus side. A different grammar would yield a different $|D_t|$. We
report ratios, so a constant rescaling of $|D_t|$ leaves them alone, but a grammar whose behavior
drifts with file age would not.

\subsection{Cross-version matching}
\label{sec:appendix:tracking}

We match instructions across consecutive versions with a three-stage cascade. We first try exact
string equality, then equality after normalization with case, whitespace and markdown punctuation
removed, then a fuzzy match above a similarity of \MatchFuzzyThreshold{}. A match at any stage
counts as \emph{survival}. We take the highest-scoring partner still available, and each
instruction matches at most once. An unmatched old instruction is a death; an unmatched new one is
a birth.

The cascade tells a deleted instruction from a reworded one, and \autoref{sec:measurement} rests on
that separation. A line-level diff scores a rewording as one deletion plus one addition, so a
line-level count of deletions is equally consistent with zero instruction deletions and with many.
The matcher therefore bounds every result we take from the corpus. We validated it once, on
\GateBPool{} hand-annotated transitions, reaching \GateBPrecision{} precision and \GateBRecall{}
recall against a floor of \GateBFloor{} that we fixed before pulling any data
(\autoref{tab:gates}). That is a small pool against \FieldMatches{} tracked transitions, and it is
the only number on the corpus side that the canonical run did not produce. Limitations says so. A
stratified re-annotation at full scale would close it.

Any error that survives makes our estimate conservative. A missed rewording invents a death at the
age where the rewording happened, and rewordings fall more often on old, heavily-edited
instructions than on young ones. Matcher misses therefore add hazard to the old-age tail, which
biases the estimated slope \emph{toward} zero. A decline measured under that bias is a lower bound
on the true decline.

\subsection{Competing risks: rewrites and migrations}
\label{sec:appendix:censoring}

\begin{table*}[t]
  \centering\small
  \begin{tabularx}{\textwidth}{@{}>{\raggedright\arraybackslash}Xrrrr@{}}
    \toprule
    \multicolumn{5}{@{}l}{\textbf{(a)} Anatomy of a rewrite commit} \\
    \midrule
    Rewrite commits & \multicolumn{4}{r@{}}{1,353} \\
    Median instructions killed / born & \multicolumn{4}{r@{}}{46 / 14} \\
    Events leaving the file empty & \multicolumn{4}{r@{}}{24.4\%} \\
    Events replacing $\ge$ half the dead text & \multicolumn{4}{r@{}}{38.3\%} \\
    Events carrying any cross-file migration & \multicolumn{4}{r@{}}{2.9\%} \\
    \midrule
    \multicolumn{5}{@{}l}{\textbf{(b)} Event study vs.\ window half-width} \\
    \midrule
    Window $w$ & $n$ & $k{=}{-}w$ & $k{=}0$ & $k{=}{+}w$ \\
    $\pm2$ & 317 & $-3\%$ & $-49\%$ & $-33\%$ \\
    $\pm3$ & 213 & $-8\%$ & $-53\%$ & $-39\%$ \\
    $\pm4$ & 158 & $-11\%$ & $-55\%$ & $-40\%$ \\
    $\pm5$ & 126 & $-16\%$ & $-61\%$ & $-39\%$ \\
    $\pm10^{\star}$ & 52 & $-29\%$ & $-71\%$ & $-58\%$ \\
    \midrule
    \multicolumn{5}{@{}l}{\textbf{(c)} Growth within inter-rewrite segments} \\
    \midrule
    Segment & $n$ & Median & Growing & Net $\Delta D$ \\
    0 & 1,597 & $\times1.33$ & 78.5\% & $+34.0$ \\
    1 & 543 & $\times1.32$ & 81.6\% & $+34.3$ \\
    2 & 180 & $\times1.37$ & 82.2\% & $+39.8$ \\
    3+ & 178 & $\times1.46$ & 86.0\% & $+42.3$ \\
    \bottomrule
  \end{tabularx}
  \caption{
    The rewrite is a mass deletion, its shape does not depend on the window it is measured in, and
    growth after one is no slower than before. \textbf{(a)} What a rewrite commit does, over all
    such commits in the corpus. \textbf{(b)} The event study of \autoref{fig:massrewrite:event}
    re-estimated at each window half-width: $D_t$ relative to the version before the file's first
    rewrite, median over the balanced panel of files observed at every offset, emptying events
    excluded. Columns are the level $w$ commits before the event, at it, and $w$ after, each
    relative to the version immediately before it; $\star$ marks the width
    \autoref{fig:massrewrite:event} draws. $n$ falls with width because a balanced panel keeps
    only files that outlive the event by the full window; the drop and the rebound do not.
    \textbf{(c)} Growth within each
    inter-rewrite segment ($\ge 2$ commits); segment $0$ runs to the first rewrite, segment $k$
    starts at the $k$-th, so its starting size is what that rewrite left behind. The rewrite
    threshold itself is held at $0.5$ throughout and is not swept (\S\ref{sec:appendix:censoring}).
  }
  \label{tab:rewrite}
\end{table*}

We want the hazard to run over one kind of event only: a maintainer deciding that an instruction is
no longer worth its place. We censor two other kinds of disappearance as competing risks. A
\emph{rewrite} is a commit at which at least \RewriteFracThreshold{} of a file's live instructions
die at once, subject to a floor of \RewriteMinDirectives{} instructions at risk, so that a
two-instruction file losing one does not qualify. A \emph{migration} is a death whose text turns up
in a sibling file at the same commit, and we match those repository-wide.

Together the two censor \FieldCensoredDeathShare{} of observed instruction deaths. Rewrite deaths
on their own outnumber the deletions we estimate the hazard from by
\FieldRewriteToDeletionRatio{} to one. That is our largest single exposure, and it reads two ways.
It threatens the estimate, because we fit the deletion hazard on a minority of deaths and a rewrite
threshold that mistook ordinary pruning for a wholesale rewrite would pull real deletions out of
the risk set. It also reports a finding: maintainers replace rather than prune, exactly as
\autoref{eq:ratchet} predicts.

\autoref{tab:rewrite} gathers the checks a reader can run without re-tracking the corpus.
Panel~(a) shows that the event really is a mass \emph{deletion}. The median event kills far more
instructions than it births, a quarter of events leave the file with none at all, and fewer than
half replace even half of the dead text. Panel~(b) re-estimates the event study at every window
half-width. A balanced panel's composition depends on that width, and without the sweep a reader
cannot separate our result from our selection. The drop and the rebound both survive it. Panel~(c)
asks whether the ratchet is self-limiting, which would show up as slower growth after later
rewrites. Growth after them is faster instead. Cross-file migration does not account for the events
either: under 3\% of rewrite commits carry any migration death at all.

We omit one check from \autoref{tab:rewrite}: a sweep of the \RewriteFracThreshold{} threshold
itself. The tracker applies that threshold internally, so varying it means re-running the pipeline
rather than re-aggregating its output, and we have no such run. Limitations names the audit.

\subsection{Growth in count versus in elaboration}
\label{sec:appendix:decomposition}

\autoref{sec:growth} splits growth into three components and reports all three. The split
falsifies our thesis, and the total does not. Suppose files grew mainly by elaborating rules they
already had. The ratchet would then be a story about verbosity, and counting instructions would
miss the point. They do not: on the same normalized-lifetime grid and the same estimator as
\autoref{fig:hero:growth}, mean instruction length relative to each file's own first version peaks
at \FieldGrowthLenPeakPct{} (over the \FieldGrowthLenFiles{} multi-version files with a defined
ratio) against the instruction count's \FieldGrowthCountPeakPct{}. Non-instruction payload grows
\FieldGrowthPayloadPct{} alongside, so the count growth is not text migrating between the
instruction and payload classes.

We chose both axes for robustness to composition. Normalized lifetime keeps every surviving file
contributing at every point rather than letting the population thin with age, and dividing by each
file's own first version removes level composition. Without that second step, a corpus in which
large files live longer would show growth that is entirely survivorship.

\subsection{The deletion hazard}
\label{sec:appendix:hazard}

\autoref{fig:hazard} plots a smoothed Nelson--Aalen estimator of deletion risk against instruction
age. At age $a$ the risk set holds every spell we observed to reach age $a$. A spell leaves that
set in one of three ways: by deletion, which is the event; by a censored competing risk; or by
reaching the end of its file's history. The third exit dominates, and it tells us little. An
instruction still alive at the last commit has a hazard that has not fired yet, which is not the
same as one that never will.

We count age in commits touching the file, not calendar days. \autoref{eq:hazard} is defined on
that clock, since $h(a)$ measures risk per review opportunity, and an instruction sitting in a
dormant repository is never offered one. We report calendar days as corroboration only, because a
decline on the calendar clock alone would also be consistent with files simply being edited less as
they age.

We kernel-smooth the rates with a bandwidth of \HazardBandwidthCommits{} commits
(\HazardBandwidthDays{} days on the calendar clock) over ages $0$ to \HazardAgeMax{}. We drop any
grid point backed by less than one expected event rather than plot it, so the curve never runs past
its own support. We take the confidence bands and the log-hazard slope from a bootstrap stratified
by repository, resampling repositories with replacement rather than spells (\HazardBoot{} draws,
seed 0). Instructions inside one repository share an author, a segmentation outcome and a
maintenance culture. A spell-level bootstrap would treat them as independent and understate the
interval severalfold. We report the median slope across draws, with the 2.5th and 97.5th
percentiles as its interval.

\subsection{The gamma frailty fit}
\label{sec:appendix:frailty}

\begin{table*}[t]
  \centering\small
  \begin{tabular}{@{}llrrrrrr@{}}
    \toprule
    Rows & Frailty & Spells & Deaths & Slope & 95\% CI & $\theta$ & Absorbed \\
    \midrule
    all spells & no frailty (baseline) & 247,694 & 28,255 & $-0.0506$ & $[-0.0524, -0.0489]$ & -- & -- \\
     & shared within repository & 247,694 & 28,255 & $-0.0481$ & $[-0.0499, -0.0464]$ & 1.44 & 4.9\% \\
     & shared within file & 247,694 & 28,255 & $-0.0485$ & $[-0.0502, -0.0467]$ & 1.49 & 4.3\% \\
    recurring text & no frailty (baseline) & 20,807 & 2,669 & $-0.0513$ & $[-0.0571, -0.0455]$ & -- & -- \\
     & shared within repository & 20,807 & 2,669 & $-0.0505$ & $[-0.0566, -0.0445]$ & 2.85 & 1.5\% \\
     & shared by instruction text & 20,807 & 2,669 & $-0.0355$ & $[-0.0414, -0.0296]$ & 3.41 & 30.8\% \\
    \bottomrule
  \end{tabular}
  \caption{
    Composition shrinks the age slope at every clustering level and never reverses it.
    Each row is one gamma-frailty specification, fit by piecewise-exponential EM
    (\S\ref{sec:appendix:frailty}). The columns give the log-hazard slope per commit with its 95\%
    CI, the fitted frailty variance $\theta$, and the share of that block's no-frailty slope which
    the clustering absorbs. \emph{Recurring text} restricts to instructions whose text appears in
    $\ge 2$ repositories, which are the rows where a frailty shared by instruction content can be
    identified; absorbed shares compare within a block, against the no-frailty fit on those same
    rows. We censor deaths administratively at 50 commits. A slope of $0$ would put deletion risk
    flat in instruction age, and a positive slope is what instruction staleness predicts. Content
    absorbs the most. Repository clustering on those same rows absorbs almost nothing, which puts
    the absorbed variance in the instruction text rather than in its author.
  }
  \label{tab:frailty}
\end{table*}


A hazard that declines across a population signals survival heterogeneity just as readily as
duration dependence. This is the content-fragility mechanism of \autoref{sec:hazard}: if some
instructions are intrinsically fragile, they die early, and the ones still alive at old ages are
the robust remainder. We test it with a \emph{gamma frailty} model, and keep that name for the
estimator throughout so it is never read as a fourth mechanism. We give each cluster a gamma
frailty $z$ with mean 1 and variance $\theta$, which makes the population hazard
$h_{\text{pop}}(a) = h_0(a) / (1 + \theta H_0(a))$. Even a perfectly flat individual baseline
$h_0$ then produces a declining $h_{\text{pop}}$, at a rate $\theta$ sets. We therefore take the
presence of frailty for granted and ask only whether it accounts for our slope.

We fit $h_0$ and $\theta$ jointly by piecewise-exponential EM on the risk table, one interval per
commit of age, with administrative censoring at \HazardAgeMax{} commits. That is a different
estimator from the bootstrap above, and we label it as such wherever both appear. The headline
slope stays the Nelson--Aalen bootstrap. \autoref{tab:frailty} carries the EM fit's own no-frailty
baseline alongside it, so every absorbed share has a stated denominator. Our test suite holds the
recovery tests, among them a check that a flat baseline at $\theta = 1$ does not read as a
decline.

We cluster at three levels. We identify repository and file frailty on every spell. We can identify
frailty shared by instruction \emph{content} only where the same normalized text recurs across at
least two repositories, so we fit it on that subsample and report the subsample's own no-frailty
baseline beside it. Content absorbs by far the most, and repository clustering on those same rows
absorbs almost nothing, which locates the absorbed variance in the instruction text rather than in
its author. The slope survives all three.

We fit one clustering level at a time, so each $\theta$ is an upper bound on the variance
attributable to that level alone. We do not fit a crossed repository $\times$ content
specification.

\subsection{The authorship-turnover interaction}
\label{sec:appendix:interaction}

\begin{table*}[t]
  \centering\small
  \begin{tabular}{@{}lrrrr@{}}
    \toprule
    Term & Coef. & SE & $z$ & 95\% CI \\
    \midrule
    age & $-0.0490$ & 0.0028 & $-17.3$ & $[-0.0546, -0.0435]$ \\
    multi-author & $-0.0237$ & 0.0178 & $-1.3$ & $[-0.0587, +0.0113]$ \\
    multi-author $\times$ age & $-0.0211$ & 0.0018 & $-11.7$ & $[-0.0246, -0.0175]$ \\
    $\log$ commits & $+0.1920$ & 0.0071 & $+26.9$ & $[+0.1780, +0.2059]$ \\
    $\log$ commits $\times$ age & $+0.0034$ & 0.0007 & $+4.8$ & $[+0.0021, +0.0048]$ \\
    \multicolumn{5}{@{}l}{\footnotesize $\theta_{\text{file}} = 1.35$; 1,837 files, 589 multi-author; 15,118 risk cells} \\
    \bottomrule
  \end{tabular}
  \caption{
    The age slope steepens where more humans edit the file, which composition cannot produce.
    Each row is a term of one piecewise-exponential hazard, fit with a file-level gamma frailty
    on top (\S\ref{sec:appendix:interaction}). The columns give the coefficient, its standard
    error, $z$, and the 95\% CI. \emph{multi-author} takes the value 1 where $\ge 2$ distinct
    human authors touched the file, with bot commits excluded. The term under test is
    \emph{multi-author $\times$ age}: imperfect recall predicts it negative, while a fixed
    per-instruction frailty predicts zero in either direction. $\log$ commits controls for file
    activity, since multi-author files are also busier. We ran this fit as exploratory work and
    pre-registered no criterion for it.
  }
  \label{tab:interaction}
\end{table*}


Content fragility and imperfect recall make the same marginal prediction. They diverge on one
conditional prediction. A fixed per-instruction frailty is a constant, so it gives no reason for
the age slope to differ between files many people edit and files one person edits. Imperfect
recall does give a reason: people hold the rationale it says decays, and people leave.

We fit one piecewise-exponential hazard with a file-level gamma frailty on top, which estimates the
interaction net of the between-file heterogeneity the previous subsection measured. We then
interact age with an indicator for whether at least two distinct human authors have touched the
file. We exclude bot commits before the count, because automated formatters and dependency bots
touch context files, carry no rationale, and would inflate the multi-author stratum with exactly
the commits the mechanism does not apply to.

Multi-author files are also busier, and file activity could drive an age slope by itself. We
therefore enter $\log$ commits both as a level and interacted with age, on the same footing as the
covariate under test. A skeptical reader should check that interaction first. We place it in
\autoref{tab:interaction} rather than in the body only because the body has room for the term
under test alone.

This analysis is exploratory. We pre-registered no criterion for it, it is one specification
rather than a specification curve, and multi-author status proxies for turnover rather than
measuring it, counting how many people have edited a file without establishing whether the person
who wrote a given instruction has left. We report it as evidence weighing against a rival account
and not as a passed gate.
